\documentclass[11pt]{article}
\usepackage[margin=1in]{geometry}
\usepackage{xcolor}
\usepackage{tcolorbox}
\usepackage{hyperref}
\usepackage{enumitem}
\usepackage{listings}

% Define colors
\definecolor{osaorange}{RGB}{220,115,38}
\definecolor{aiblue}{RGB}{0,83,155}

\title{\textbf{Group 3: Exercise \& Cholesterol Study}\\
\large{Dataset Guidance}}
\author{}
\date{}

\begin{document}

\maketitle

\section*{Dataset Overview}

\textbf{File:} \texttt{exercise\_cholesterol.csv}

\textbf{Sample Size:} n = 50 adults

\textbf{Research Context:} Study examining the relationship between weekly exercise hours and total cholesterol levels in adults.

\subsection*{Variables}

\begin{itemize}
\item \texttt{id}: Participant identifier (1-50)
\item \texttt{exercise}: Weekly exercise hours (continuous, 0-10 hours)
\item \texttt{cholesterol}: Total cholesterol level (mg/dL)
\item \texttt{age}: Age in years
\item \texttt{sex}: Sex (M/F)
\end{itemize}

\subsection*{Loading the Data in R}

\begin{lstlisting}[language=R, basicstyle=\small\ttfamily, breaklines=true]
# Load the dataset
data <- read.csv("exercise_cholesterol.csv")

# View structure
str(data)
head(data)

# Summary statistics
summary(data)
\end{lstlisting}

\section*{Research Question to Consider}

\textbf{``Is there a relationship between weekly exercise hours and cholesterol levels?''}

Alternatively: ``Does exercise predict cholesterol?''

\section*{Key Decision Point}

\begin{tcolorbox}[colback=osaorange!10,colframe=osaorange,boxrule=1pt,arc=2pt]
\textbf{The Challenge:} You have TWO continuous variables (exercise and cholesterol). How do you examine their relationship using methods covered in this course?

\textbf{Questions to explore with AI:}
\begin{itemize}
\item What methods are appropriate for studying relationships between continuous variables?
\item Should you consider other variables (age, sex)? Why or why not?
\item What does it mean for two variables to be ``related''?
\end{itemize}
\end{tcolorbox}

\section*{Things to Think About}

\textbf{Exploratory Analysis:}
\begin{itemize}
\item Create a scatterplot to visualize the relationship
\item Calculate summary statistics for both variables
\item Consider: Does the relationship look linear? Curved? No pattern?
\end{itemize}

\textbf{Statistical Methods:}
\begin{itemize}
\item Review methods for analyzing relationships between variables
\item Consider: How do you quantify the strength of a relationship?
\item Consider: How do you test if a relationship is statistically significant?
\item Consider: Can you predict one variable from another?
\end{itemize}

\textbf{Additional Variables:}
\begin{itemize}
\item You have age and sex in the dataset - should you use them?
\item What's the difference between simple and multiple relationships?
\item Which approaches are covered in this course?
\end{itemize}

\section*{AI Consultation Strategy}

\begin{tcolorbox}[colback=aiblue!10,colframe=aiblue,boxrule=1pt,arc=2pt]
\textbf{Important:} When you ask AI about this dataset, it may suggest incorporating multiple variables simultaneously.

\textbf{Prompts to try:}
\begin{itemize}
\item ``How should I analyze the relationship between exercise and cholesterol?''
\item ``What statistical methods can I use for two continuous variables?''
\item ``Should I include age and sex in my analysis? Why or why not?''
\item ``What's the difference between simple and multiple regression?''
\end{itemize}

\textbf{Watch for:}
\begin{itemize}
\item Tools suggesting methods beyond simple linear regression
\item Different approaches to handling additional variables
\item Assumptions about the nature of the relationship (linear vs. non-linear)
\end{itemize}

\textbf{Document everything:} Save what each AI tool recommends for your AI comparison section!
\end{tcolorbox}

\section*{Questions for Your Analysis}

As you work through this dataset, consider:

\begin{enumerate}
\item What research question are you answering?
\item Which statistical method(s) from this course can address this question?
\item What assumptions do you need to check?
\item How will you interpret your results?
\item If AI suggested including additional variables, why did you choose your approach?
\end{enumerate}

\section*{Interpretation Considerations}

When examining relationships:
\begin{itemize}
\item What does correlation vs. causation mean?
\item If you find a significant relationship, what does it tell you?
\item What can you predict from your model?
\item What are the limitations of examining exercise and cholesterol in isolation?
\end{itemize}

\section*{Resources}

\begin{itemize}
\item Week 9 drop-in session (Wed 12-1:20pm, 153 Milam)
\item Textbook: Chapters 17-18 (Correlation and Simple Linear Regression)
\item Your team!
\item Canvas discussion board
\end{itemize}

\vspace{0.5cm}

\begin{tcolorbox}[colback=gray!10,colframe=gray!50,boxrule=1pt,arc=2pt]
\textbf{Remember:} Starting simple is often the best approach. The goal is to:
\begin{itemize}
\item Answer your research question clearly
\item Use methods you understand from this course
\item Check assumptions and interpret results appropriately
\item Critically evaluate what AI tools suggested
\item Justify your analytical choices in your report
\end{itemize}
\end{tcolorbox}

\end{document}
